\documentclass[11pt]{article}

\usepackage[margin=1in]{geometry}
\usepackage{amsmath,amssymb,amsthm,mathtools}
\usepackage{hyperref}
\usepackage{enumitem}
\usepackage{microtype}

\raggedbottom
\allowdisplaybreaks[2]

\setlength{\parindent}{1.5em}
\setlength{\parskip}{0pt}
\setlength{\abovedisplayskip}{7pt plus 2pt minus 3pt}
\setlength{\belowdisplayskip}{7pt plus 2pt minus 3pt}
\setlength{\abovedisplayshortskip}{4pt plus 2pt minus 2pt}
\setlength{\belowdisplayshortskip}{4pt plus 2pt minus 2pt}

\setlist{topsep=3pt,itemsep=2pt,parsep=0pt,partopsep=0pt}

\hypersetup{
  colorlinks=true,
  linkcolor=blue,
  citecolor=blue,
  urlcolor=blue
}

\newtheoremstyle{tightplain}
  {6pt}{6pt}{\itshape}{}{\bfseries}{.}{.5em}{}
\newtheoremstyle{tightdef}
  {6pt}{6pt}{}{}{\bfseries}{.}{.5em}{}

\theoremstyle{tightplain}
\newtheorem{theorem}{Theorem}[section]
\newtheorem{proposition}[theorem]{Proposition}
\newtheorem{lemma}[theorem]{Lemma}
\newtheorem{corollary}[theorem]{Corollary}

\theoremstyle{tightdef}
\newtheorem{assumption}[theorem]{Assumption}
\newtheorem{remark}[theorem]{Remark}
\newtheorem{definition}[theorem]{Definition}

\DeclareMathOperator{\tr}{tr}
\DeclareMathOperator{\diag}{diag}
\DeclareMathOperator{\rank}{rank}
\DeclareMathOperator{\op}{op}
\DeclareMathOperator{\KL}{KL}

\newcommand{\E}{\mathbb E}
\newcommand{\Prob}{\mathbb P}
\newcommand{\R}{\mathbb R}
\newcommand{\norm}[1]{\left\lVert #1 \right\rVert}
\newcommand{\abs}[1]{\left\lvert #1 \right\rvert}
\newcommand{\one}{\mathbf 1}
\newcommand{\F}{\mathrm F}
\newcommand{\Shat}{\widehat\Sigma}
\newcommand{\Phat}{\widehat P}
\newcommand{\Sbar}{\bar\Sigma}
\newcommand{\dsub}{\mathfrak d_{\mathrm{sub}}}
\newcommand{\ssum}{\mathsf s}

\title{Overlap-Driven Subspace Drift in Rolling Covariance Windows}
\author{Anonymous}
\date{\today}

\begin{document}

\maketitle

\begin{abstract}
We analyze the increments of top-$r$ spectral projectors of rolling sample
covariance matrices. Two adjacent windows of length $W$ share $W-1$
observations and differ only through a rank-two boundary innovation, so the
difference of adjacent rolling sample covariances has rank at most two.
Working under a static Gaussian bulk-plus-spikes null, we use this shared-core
structure to obtain four results: a conditional exchangeability principle for
the two adjacent projectors; an exact finite-sample sign-randomization
property for oriented projector contrasts; an adjacent-window drift benchmark
of order $r/W^2$; and the leading-order expansion
\[
\E\norm{\widehat P_t-\widehat P_{t-1}}_F^2
=
\frac{2\mathfrak d_{\mathrm{sub}}(\Sigma,r)}{W^2}+o(W^{-2}),
\]
where $\mathfrak d_{\mathrm{sub}}$ is the spectral difficulty of
Definition~\ref{def:dsub}. A boundary-driven minimax lower bound shows that the
rate $r/W^2$ is sharp for the rolling drift target
$P(\bar\Sigma_t)-P(\bar\Sigma_{t-1})$. The mechanism throughout is to condition
on the shared core of $W-1$ common observations, which cancels in the adjacent
difference and reduces the effective noise to the two boundary observations.
\end{abstract}

\section{Introduction}

\subsection{The adjacent subspace increment}

Let $\{r_s\}$ be a sequence of mean-zero random vectors in $\R^n$, let
$I_t=\{t-W+1,\ldots,t\}$ denote the rolling window of length $W$ ending at time
$t$, and define the rolling sample covariance
\[
\widehat\Sigma_t=\frac1W\sum_{s\in I_t}r_s r_s^\top.
\]
For a symmetric matrix $A$ with a positive top-$r$ eigengap, let $P(A)$ denote
its top-$r$ spectral projector, and write $\widehat P_t=P(\widehat\Sigma_t)$.
The object of this paper is the \emph{adjacent subspace increment}
\[
\widehat P_t-\widehat P_{t-1}.
\]
Our focus is neither the covariance estimate itself nor single-window principal
component analysis, but the increment of the estimated principal subspace
between two consecutive windows. The windows $I_{t-1}$ and $I_t$ overlap in
$W-1$ observations, so that
\begin{equation}\label{eq:rank2}
\widehat\Sigma_t-\widehat\Sigma_{t-1}
=\frac1W\bigl(r_t r_t^\top-r_{t-W}r_{t-W}^\top\bigr).
\end{equation}
Although each estimate $\widehat\Sigma_t$ is built from $W$ observations, their
difference has rank at most two and depends only on the entering observation
$r_t$ and the exiting observation $r_{t-W}$. The rank-two boundary identity
\eqref{eq:rank2} is the source of the faster decay rate that we establish for
the adjacent increment.

\subsection{Related work}

Principal subspaces and spectral projectors of covariance matrices are central
objects in multivariate statistics. Their behavior in high dimensions has been
studied extensively under the spiked covariance model of Johnstone
\cite{Johnstone2001}; see Paul \cite{Paul2007} and Benaych-Georges and
Nadakuditi \cite{BenaychNadakuditi2011} for the eigenstructure of low-rank
perturbations of large random matrices. Finite-sample and minimax theory for
principal subspace estimation includes Vu and Lei \cite{VuLei2013}, Cai, Ma and
Wu \cite{CaiMaWu2013}, and Cai and Zhang \cite{CaiZhang2018}. The underlying
perturbation theory for spectral projectors goes back to Davis and Kahan
\cite{DavisKahan1970}, Wedin \cite{Wedin1972}, and Kato \cite{Kato1995}; see
also Stewart and Sun \cite{StewartSun1990} and Bhatia \cite{Bhatia1997}. Sharp
operator-norm concentration for sample covariance operators in terms of
effective rank is due to Koltchinskii and Lounici \cite{KoltchinskiiLounici2017}.

Rolling-window covariance and subspace estimates are widely used to track
time-varying second-order structure and to localize changes in covariance
\cite{AueEtAl2009,WangYuRinaldo2021}. In such analyses the temporal overlap
between consecutive windows is typically not exploited. The present paper
isolates the effect of this overlap on the projector increment and shows that,
under a static null, it produces a drift scale that is one order of magnitude
in $W$ smaller than the single-window estimation error.

\subsection{Contributions}

We work under a static Gaussian bulk-plus-spikes null, made precise in
Section~\ref{sec:setup}. Under this null the population subspace does not move,
so the increment $\widehat P_t-\widehat P_{t-1}$ is pure estimation noise, and
its size is the natural benchmark against which true subspace drift must be
detected. Our contributions are the following.
\begin{enumerate}
\item \emph{Conditional exchangeability}
(Proposition~\ref{prop:exchangeability}). Conditional on the shared core of
$W-1$ common observations, the two adjacent projectors are independent and
identically distributed; in particular the increment is conditionally
mean-zero.
\item \emph{Exact finite-sample sign randomization}
(Theorem~\ref{thm:exactrandomization} and Corollary~\ref{cor:signtest}). Any
antisymmetric functional of the two adjacent projectors --- in particular any
linear contrast of $\widehat P_t-\widehat P_{t-1}$ --- has a sign-symmetric
conditional law, giving an exact finite-sample sign test for oriented
contrasts.
\item \emph{Adjacent-window drift benchmark}
(Theorem~\ref{thm:overlapbenchmark}). The expected squared Frobenius increment
is of order $r/W^2$, in contrast to the single-window projector error of order
$r/W$.
\item \emph{Constant-level expansion}
(Theorem~\ref{thm:adjacentconstantproved}). The leading constant is identified
exactly,
\[
\E\norm{\widehat P_t-\widehat P_{t-1}}_F^2
= \frac{2\,\mathfrak d_{\mathrm{sub}}(\Sigma,r)}{W^2}+o(W^{-2}),
\]
where $\mathfrak d_{\mathrm{sub}}$ is the linearized spectral difficulty of
Definition~\ref{def:dsub}.
\item \emph{Boundary-driven lower bound}
(Proposition~\ref{prop:rollinglower}). The rate $r/W^2$ is minimax sharp for
the rolling drift target $D_t^{\mathrm{roll}}=P(\bar\Sigma_t)-P(\bar\Sigma_{t-1})$.
\end{enumerate}
Standard perturbation and concentration results used as inputs are collected in
Section~\ref{sec:aux}.

\subsection{Notation}

For a symmetric matrix $A$, $\lambda_1(A)\ge\cdots\ge\lambda_n(A)$ are its
ordered eigenvalues, $\norm{A}_{\op}$ its operator norm, $\norm{A}_F$ its
Frobenius norm, and $\langle A,B\rangle_F=\tr(A^\top B)$ the Frobenius inner
product. We write $a\lesssim b$ if $a\le Cb$ for a constant $C$ depending only
on the fixed spectral constants and the rank $r$, $a\asymp b$ if $a\lesssim b$
and $b\lesssim a$, and $a=\Theta(b)$ analogously. For a Hilbert-space-valued
random variable $X$ and a sub-$\sigma$-algebra $\mathcal G$,
\[
V_F(X\mid\mathcal G)
=
\E\bigl[\norm{X-\E[X\mid\mathcal G]}_F^2\bigm|\mathcal G\bigr]
\]
is the conditional Frobenius variance.

\section{Setup}\label{sec:setup}

Throughout the static-null theory we assume that on $I_t\cup I_{t-1}$,
\[
r_s\stackrel{\mathrm{ind}}{\sim}N(0,\Sigma).
\]
Let $P(A)$ denote the top-$r$ spectral projector of a symmetric matrix $A$
whenever it is uniquely defined. On null sets with eigenvalue ties, fix any
measurable tie-breaking convention.

\begin{assumption}[Static Gaussian null]\label{ass:gaussian}
On $I_t\cup I_{t-1}$, the vectors $r_s$ are independent and identically
distributed as $N(0,\Sigma)$.
\end{assumption}

\begin{assumption}[Bulk-plus-spikes spectrum]\label{ass:bulkspike}
The rank parameter $r$ is fixed. Eigenvalues are ordered as
\[
\lambda_1(\Sigma)\ge\lambda_2(\Sigma)\ge\cdots\ge\lambda_n(\Sigma).
\]
There exist constants $0<m<M<\infty$ and $0<c_\lambda\le C_\lambda<\infty$,
independent of $n$ and $W$, such that
\[
m\le \lambda_n(\Sigma)\le\cdots\le\lambda_{r+1}(\Sigma)\le M,
\qquad
c_\lambda n\le \lambda_r(\Sigma)\le\cdots\le\lambda_1(\Sigma)\le C_\lambda n.
\]
\end{assumption}

\begin{assumption}[Population eigengap]\label{ass:gap}
Let $\delta=\lambda_r(\Sigma)-\lambda_{r+1}(\Sigma)$. Assume
\[
\delta\ge c_\delta n.
\]
\end{assumption}

\begin{assumption}[Asymptotic convention]\label{ass:asymptotic}
The spike rank $r$ is fixed, $n$ may grow, and $W=W_n\to\infty$. Whenever
exponential sample-covariance concentration is invoked, we assume
$W\ge C\log n$ and $W\ge W_0$, with constants depending only on the fixed
spectral constants and on $r$.
\end{assumption}

The bulk-plus-spikes scaling of Assumption~\ref{ass:bulkspike} places the model
in the spiked regime studied by Johnstone \cite{Johnstone2001} and Paul
\cite{Paul2007}: $r$ spikes of order $n$ separated from an $O(1)$ bulk by a gap
of order $n$.

\section{Auxiliary standard inputs}\label{sec:aux}

We collect the standard perturbation and concentration results used in the
sequel. They are stated in the form convenient for our purposes and are
included to keep the paper self-contained.

\begin{lemma}[First-order projector expansion]\label{lem:projectorexpansion}
Let $A=A^\top$ have eigen-decomposition
\[
A=\sum_{k=1}^n \mu_k v_kv_k^\top,
\qquad
\mu_1\ge\cdots\ge\mu_n,
\]
and suppose the top-$r$ eigengap $\gamma=\mu_r-\mu_{r+1}$ is positive. Let
$P(A)$ be the top-$r$ spectral projector of $A$. Define the reduced-resolvent
derivative
\[
\mathcal L_A(E)
=\sum_{i\le r<j}
\frac{v_j v_j^\top E\, v_i v_i^\top+v_i v_i^\top E\, v_j v_j^\top}{\mu_i-\mu_j}.
\]
There exist constants $c_0\in(0,1/4]$ and $C_r<\infty$, depending at most on the
fixed rank $r$, such that if $\norm{E}_{\op}\le c_0\gamma$, then the top-$r$
projector $P(A+E)$ is uniquely defined and
\[
\norm{P(A+E)-P(A)-\mathcal L_A(E)}_F
\le
C_r\frac{\norm{E}_{\op}^2}{\gamma^2}.
\]
Consequently,
\[
\norm{P(A+E)-P(A)}_F
\le
C_r\frac{\norm{E}_{\op}}{\gamma}.
\]
\end{lemma}

\begin{proof}
Let $P=P(A)$ and write $\widetilde P=P(A+E)$. By Weyl's inequality
\cite[Ch.~III]{Bhatia1997}, if $\norm E_{\op}\le \gamma/4$ then the top-$r$
spectral cluster of $A+E$ remains separated from the lower cluster by at least
$\gamma/2$, so $\widetilde P$ is uniquely defined. We use the standard local
perturbation expansion for spectral projectors across an isolated gap
\cite{Kato1995,DavisKahan1970}: the map $H\mapsto P(A+H)$ is twice Fr\'echet
differentiable on the operator-norm ball $\norm H_{\op}\le\gamma/4$, its
derivative at zero is the reduced-resolvent operator $\mathcal L_A$, and
\[
\norm{D^2P(A+H)[B_1,B_2]}_{\op}
\le
C\frac{\norm{B_1}_{\op}\norm{B_2}_{\op}}{\gamma^2}
\]
for all $\norm H_{\op}\le\gamma/4$. This is the standard second-order bound for
spectral projectors across an isolated spectral gap; equivalently, it follows
from the Riesz-projector (contour-integral) representation
\cite[Ch.~II]{Kato1995} using any contour that strictly separates the top-$r$
cluster from its complement.

Taylor's formula therefore gives
\[
\norm{P(A+E)-P(A)-\mathcal L_A(E)}_{\op}
\le
C\frac{\norm E_{\op}^2}{\gamma^2}.
\]
The linear term has the displayed coordinate form by differentiating the
spectral projector in the eigenbasis of $A$. Finally, $P(A+E)-P(A)$ is the
difference of two rank-$r$ projectors and has rank at most $2r$, while
$\mathcal L_A(E)$ has only the two off-diagonal blocks between $P$ and
$P^\perp$ and rank at most $2r$. Hence the remainder has rank at most $4r$, so
\[
\norm{P(A+E)-P(A)-\mathcal L_A(E)}_F
\le
2\sqrt r\,\norm{P(A+E)-P(A)-\mathcal L_A(E)}_{\op}
\le
C_r\frac{\norm E_{\op}^2}{\gamma^2}.
\]
The final Lipschitz bound follows from
$\norm{\mathcal L_A(E)}_F\le C_r\norm E_{\op}/\gamma$ and the remainder bound,
after decreasing $c_0$ if necessary.
\end{proof}

\begin{lemma}[Gaussian operator-moment tools]\label{lem:covopmoments}
Let $S_N=N^{-1}\sum_{k=1}^N x_kx_k^\top$ with
$x_k\stackrel{\mathrm{iid}}{\sim}N(0,\Sigma)$, and suppose $\Sigma$ satisfies
Assumption~\ref{ass:bulkspike} with fixed $r$. Then, for all $N\ge N_0$,
\[
\E\norm{S_N-\Sigma}_{\op}^4\le C\frac{n^4}{N^2}.
\]
In particular $\E\norm{S_N-\Sigma}_{\op}^2\le Cn^2/N$ by Jensen's inequality.
Moreover, for every sufficiently small fixed $\eta>0$ there are constants
$c,C>0$ such that
\[
\Prob\bigl\{\norm{S_N-\Sigma}_{\op}>\eta n\bigr\}\le Ce^{-cN}
\]
whenever $N\ge C\log n$ and $N\ge N_0$.
\end{lemma}

\begin{proof}
The effective rank satisfies
\[
r_{\mathrm{eff}}(\Sigma)=\frac{\tr\Sigma}{\norm{\Sigma}_{\op}}=O(1)
\]
under Assumption~\ref{ass:bulkspike}, since $\tr\Sigma\asymp n$ and
$\norm{\Sigma}_{\op}\asymp n$. The Gaussian effective-rank covariance
concentration inequality \cite[Thm.~9]{KoltchinskiiLounici2017} (see also
\cite[Ch.~4--5]{Vershynin2018}) gives, with probability at least $1-e^{-u}$,
\[
\norm{S_N-\Sigma}_{\op}
\le
C\norm{\Sigma}_{\op}
\left(
\sqrt{\frac{r_{\mathrm{eff}}(\Sigma)+u}{N}}
+
\frac{r_{\mathrm{eff}}(\Sigma)+u}{N}
\right).
\]
Since $\norm{\Sigma}_{\op}\asymp n$, integrating this tail bound gives the
fourth-moment estimate, and the second-moment bound follows from
$\E\norm{S_N-\Sigma}_{\op}^2\le(\E\norm{S_N-\Sigma}_{\op}^4)^{1/2}$. Taking
$u=c_0N$ with $c_0>0$ sufficiently small gives the exponential deviation bound
once $N\ge N_0$; the additional condition $N\ge C\log n$ matches the
ambient-dimensional formulations of the same inequality that carry an explicit
dimension-dependent prefactor.
\end{proof}

\section{Linearized spectral difficulty}

\begin{definition}[Spectral difficulty]\label{def:dsub}
Let $\Sigma=\sum_{k=1}^n\lambda_k u_k u_k^\top$ with
$\lambda_1\ge\cdots\ge\lambda_n$ and $\lambda_r>\lambda_{r+1}$. Define
\[
\mathfrak d_{\mathrm{sub}}(\Sigma,r)
=2\sum_{i\le r<j}\frac{\lambda_i\lambda_j}{(\lambda_i-\lambda_j)^2}.
\]
\end{definition}

\begin{remark}[Factor-of-two convention]\label{rem:convention}
The leading factor $2$ is kept \emph{inside} the definition so that the
single-window linearized projector risk is exactly
$\mathfrak d_{\mathrm{sub}}(\Sigma,r)/W$
(Proposition~\ref{prop:dsub}). Write
\[
\ssum(\Sigma,r):=\sum_{i\le r<j}\frac{\lambda_i\lambda_j}{(\lambda_i-\lambda_j)^2}
=\tfrac12\,\mathfrak d_{\mathrm{sub}}(\Sigma,r)
\]
for the un-doubled cross-spectral sum. The adjacent drift constant of
Theorem~\ref{thm:adjacentconstantproved} is then
$2\mathfrak d_{\mathrm{sub}}/W^2=4\ssum/W^2$ for the full squared increment, and
$\mathfrak d_{\mathrm{sub}}/W^2=2\ssum/W^2$ for the half-norm increment
$\tfrac12\norm{\widehat P_t-\widehat P_{t-1}}_F^2$. Applications that adopt the
un-doubled sum $\ssum$ as their working spectral-difficulty functional must
carry the compensating factor $2$ in these benchmarks; mismatching the
convention rescales every null benchmark by a factor of two.
\end{remark}

\begin{proposition}[Exact linearized projector risk]\label{prop:dsub}
Let $P$ be the top-$r$ projector of a fixed covariance $\Sigma$ and let
\[
\mathcal L_\Sigma(E)=
\sum_{i\le r<j}
\frac{u_j u_j^\top E\, u_i u_i^\top+u_i u_i^\top E\, u_j u_j^\top}{\lambda_i-\lambda_j}.
\]
Then
\[
\norm{\mathcal L_\Sigma(E)}_F^2
=2\sum_{i\le r<j}\frac{\abs{u_j^\top E u_i}^2}{(\lambda_i-\lambda_j)^2}.
\]
If $\widehat\Sigma=W^{-1}\sum_{s=1}^W r_s r_s^\top$ with i.i.d.\
$r_s\sim N(0,\Sigma)$, then
\[
\E\norm{\mathcal L_\Sigma(\widehat\Sigma-\Sigma)}_F^2
=\frac{\mathfrak d_{\mathrm{sub}}(\Sigma,r)}{W}.
\]
If
\[
H=\frac1W\bigl(r_+r_+^\top-r_-r_-^\top\bigr),
\qquad
r_+,r_-\stackrel{\mathrm{iid}}{\sim}N(0,\Sigma),
\]
then
\[
\E\norm{\mathcal L_\Sigma(H)}_F^2
=\frac{2\,\mathfrak d_{\mathrm{sub}}(\Sigma,r)}{W^2}.
\]
\end{proposition}

\begin{proof}
The Frobenius identity follows from the orthonormality of the matrix units
$\{u_j u_i^\top,\,u_i u_j^\top:i\le r<j\}$ in the reduced-resolvent expression:
each pair $(i,j)$ contributes the coefficient $(u_j^\top E u_i)/(\lambda_i-\lambda_j)$
on both $u_j u_i^\top$ and $u_i u_j^\top$ (equal because $E$ is symmetric),
giving the displayed sum. For Gaussian data and $i\ne j$, Isserlis' (Wick's)
formula \cite{Isserlis1918,Janson1997} together with the orthogonality
$u_i^\top\Sigma u_j=\lambda_i\delta_{ij}$ gives
\[
\E\abs{u_j^\top(\widehat\Sigma-\Sigma)u_i}^2
=\frac{\lambda_i\lambda_j}{W},
\qquad
\E\abs{u_j^\top H u_i}^2
=\frac{2\lambda_i\lambda_j}{W^2}.
\]
Summing over $i\le r<j$ and using the definition of
$\mathfrak d_{\mathrm{sub}}$ proves both claims.
\end{proof}

\begin{corollary}[Bulk-plus-spikes first-order scale]\label{cor:firstorderscale}
Under Assumptions~\ref{ass:bulkspike}--\ref{ass:gap},
\[
\mathfrak d_{\mathrm{sub}}(\Sigma,r)=\Theta(r).
\]
Hence the first-order single-window linearized projector scale is $r/W$, while
the first-order adjacent-window linearized scale is $r/W^2$.
\end{corollary}

\begin{proof}
For $i\le r<j$,
\[
\lambda_i\asymp n,
\qquad
\lambda_j\asymp 1,
\qquad
\lambda_i-\lambda_j\asymp n.
\]
Thus each summand is $\asymp 1/n$. There are $r(n-r)$ summands, giving
$\Theta(r)$ for fixed $r$.
\end{proof}

\section{Adjacent-window overlap theory}

\subsection{Rank-two identity and shared core}

For adjacent windows, identity \eqref{eq:rank2} reads
\[
\widehat\Sigma_t-\widehat\Sigma_{t-1}
=\frac1W\bigl(r_t r_t^\top-r_{t-W}r_{t-W}^\top\bigr),
\]
so the covariance difference has rank at most two.

Under the static null $\Sigma_s\equiv\Sigma$ on $I_t\cup I_{t-1}$, define the
shared core and its normalization constant by
\[
C_t:=\frac1W\sum_{s=t-W+1}^{t-1}r_s r_s^\top,
\qquad
\alpha_W:=\frac{W-1}{W}.
\]
Then $\E C_t=\alpha_W\Sigma$, not $\Sigma$. This centering is important because
\[
\widehat\Sigma_t=C_t+\frac1W r_t r_t^\top,
\qquad
\widehat\Sigma_{t-1}=C_t+\frac1W r_{t-W}r_{t-W}^\top.
\]
Let $\phi_{C_t}(v)$ denote the top-$r$ spectral projector of $C_t+W^{-1}vv^\top$,
with the same measurable tie-breaking convention if needed.

\begin{proposition}[Conditional exchangeability]\label{prop:exchangeability}
Under the static null, conditional on
\[
\mathcal G_t=\sigma(r_{t-W+1},\dots,r_{t-1}),
\]
the projectors
\[
\widehat P_t=\phi_{C_t}(r_t),
\qquad
\widehat P_{t-1}=\phi_{C_t}(r_{t-W})
\]
are independent and identically distributed. Consequently,
\[
\E[\widehat P_t-\widehat P_{t-1}\mid\mathcal G_t]=0,
\qquad
\E\norm{\widehat P_t-\widehat P_{t-1}}_F^2
=2\,\E\bigl[V_F(\widehat P_t\mid\mathcal G_t)\bigr].
\]
\end{proposition}

\begin{proof}
Conditional on $\mathcal G_t$, the core $C_t$ is deterministic and
$r_t,r_{t-W}$ are i.i.d.\ $N(0,\Sigma)$, independent of the core; hence
$\widehat P_t$ and $\widehat P_{t-1}$ are i.i.d.\ given $\mathcal G_t$. For two
conditionally i.i.d.\ square-integrable Hilbert-space-valued variables
$X_1,X_2$ one has
$\E[\norm{X_1-X_2}_F^2\mid\mathcal G_t]=2\,V_F(X_1\mid\mathcal G_t)$, because
$\E[\langle X_1,X_2\rangle_F\mid\mathcal G_t]=\norm{\E[X_1\mid\mathcal G_t]}_F^2$;
taking the unconditional expectation gives the stated identity.
\end{proof}

\begin{theorem}[Exact conditional sign randomization]\label{thm:exactrandomization}
Under the static null, let $\mathcal A$ be any measurable Hilbert-space-valued
functional of two rank-$r$ projectors satisfying
\[
\mathcal A(Q_1,Q_2)=-\mathcal A(Q_2,Q_1).
\]
Then, conditional on $\mathcal G_t$,
\[
\mathcal A(\widehat P_t,\widehat P_{t-1})
\stackrel{d}=
-\mathcal A(\widehat P_t,\widehat P_{t-1}).
\]
This applies to oriented statistics such as $\widehat P_t-\widehat P_{t-1}$ and
to linear contrasts of this difference. It does not by itself calibrate the
nonnegative magnitude $\norm{\widehat P_t-\widehat P_{t-1}}_F$, nor does it give
the full finite-sample null law of such norms.
\end{theorem}

\begin{proof}
Swapping the two conditionally i.i.d.\ boundary vectors $r_t$ and $r_{t-W}$
preserves the conditional law and reverses the order of the two projectors.
Antisymmetry of $\mathcal A$ negates the statistic.
\end{proof}

\begin{corollary}[Finite-sample sign test for oriented contrasts]\label{cor:signtest}
Under the static null, fix any deterministic symmetric matrix $A$ or any
$\mathcal G_t$-measurable symmetric matrix $A_t$, and define
\[
T_t=\langle A_t,\widehat P_t-\widehat P_{t-1}\rangle_F.
\]
Then, conditional on $\mathcal G_t$,
\[
T_t\stackrel d=-T_t.
\]
Consequently, if $\Prob(T_t=0\mid\mathcal G_t)=0$, the conditional signs of
$T_t$ are exactly balanced under the static null, which yields an exact
finite-sample sign test in the sense of randomization tests
\cite[Ch.~15]{LehmannRomano2005}. If there is an atom at zero, exact
calibration requires an explicit zero convention, such as randomizing the sign
when $T_t=0$; treating zeros as non-rejections gives a conservative test. The
measurability requirement on $A_t$ is essential: a contrast that depends on the
boundary swap (for instance, the leading eigenvector of
$\widehat P_t-\widehat P_{t-1}$) co-rotates with the difference and the symmetry
fails. The statement concerns oriented contrasts and does not give the full
null law of the nonnegative norm $\norm{\widehat P_t-\widehat P_{t-1}}_F$.
\end{corollary}

\begin{proof}
Apply Theorem~\ref{thm:exactrandomization} to the antisymmetric functional
$\mathcal A(Q_1,Q_2)=\langle A_t,Q_1-Q_2\rangle_F$, conditioning on
$\mathcal G_t$ when $A_t$ is random; the swap leaves $A_t$ fixed precisely
because $A_t$ is $\mathcal G_t$-measurable.
\end{proof}

\subsection{Shared-core separation}

Let
\[
C_t=\sum_{k=1}^n\widetilde\lambda_k\widetilde u_k\widetilde u_k^\top,
\qquad
\widetilde\lambda_1\ge\cdots\ge\widetilde\lambda_n,
\]
and define the good event
\[
\mathcal E_t^{\mathrm{core}}
=\left\{\norm{C_t-\alpha_W\Sigma}_{\op}\le \eta\,\alpha_W n\right\}.
\]
Under the static null $\delta_t:=\lambda_r(\Sigma)-\lambda_{r+1}(\Sigma)=\delta\ge c_\delta n$.
Choose $\eta>0$ sufficiently small relative to $c_\lambda$ and $c_\delta$; for
example, it is enough to take
\[
\eta\le\frac{c_\delta}{4},
\qquad
\eta\le\frac{c_\lambda}{4}.
\]
On $\mathcal E_t^{\mathrm{core}}$, Weyl's inequality \cite[Ch.~III]{Bhatia1997}
gives
\[
\widetilde\lambda_r-\widetilde\lambda_{r+1}
\ge
\alpha_W\delta_t-2\eta\,\alpha_W n
\ge
\frac{\alpha_W\delta_t}{2},
\]
where the last step uses $\eta\le c_\delta/4\le \delta_t/(4n)$. Since
$W\to\infty$, the factor $\alpha_W=(W-1)/W$ is bounded away from zero and does
not affect the order of the bounds below. Define the top-$r$ shared-core
projector
\[
\widetilde P_t^C
=\sum_{i=1}^r \widetilde u_i\widetilde u_i^\top
\]
whenever the shared-core eigengap is positive.

\begin{lemma}[Shared-core separation probability]\label{lem:sharedcoreevent}
Under the static null and Assumptions~\ref{ass:bulkspike}--\ref{ass:gap}, there
exist constants $c,C>0$ such that
\[
\Prob\bigl((\mathcal E_t^{\mathrm{core}})^c\bigr)\le Ce^{-cW}
\]
whenever $W\ge C\log n$ and $W\ge W_0$ for a sufficiently large constant $W_0$.
\end{lemma}

\begin{proof}
Write
\[
C_t=\alpha_W S_t^{\mathrm{core}},
\qquad
S_t^{\mathrm{core}}=\frac1{W-1}\sum_{s=t-W+1}^{t-1}r_s r_s^\top,
\]
so that
$\norm{C_t-\alpha_W\Sigma}_{\op}=\alpha_W\norm{S_t^{\mathrm{core}}-\Sigma}_{\op}$.
The core consists of $W-1$ i.i.d.\ $N(0,\Sigma)$ vectors, and
Lemma~\ref{lem:covopmoments} (with $N=W-1$) gives
$\Prob\{\norm{S_t^{\mathrm{core}}-\Sigma}_{\op}>\eta n\}\le Ce^{-cW}$ once
$W\ge C\log n$ and $W\ge W_0$. Therefore
$\norm{C_t-\alpha_W\Sigma}_{\op}\le\eta\,\alpha_W n$ with probability at least
$1-Ce^{-cW}$.
\end{proof}

\subsection{Conditional linearized risk}

On the event $\widetilde\lambda_r>\widetilde\lambda_{r+1}$, let
$\mathcal L_{C_t}$ be the Fr\'echet derivative of the top-$r$ projector map at
$C_t$, expressed in the \emph{shared-core} eigenbasis:
\[
\mathcal L_{C_t}(E)
=\sum_{i\le r<j}
\frac{
\widetilde u_j\widetilde u_j^\top E\,\widetilde u_i\widetilde u_i^\top
+
\widetilde u_i\widetilde u_i^\top E\,\widetilde u_j\widetilde u_j^\top
}
{\widetilde\lambda_i-\widetilde\lambda_j}.
\]
Define, also in the shared-core eigenbasis,
\[
a_i=\widetilde u_i^\top\Sigma\widetilde u_i,
\qquad
b_j=\widetilde u_j^\top\Sigma\widetilde u_j,
\qquad
c_{ij}=\widetilde u_i^\top\Sigma\widetilde u_j.
\]
Note that $c_{ij}$ need not vanish, since the shared-core eigenvectors do not in
general diagonalize $\Sigma$.

\begin{proposition}[Exact conditional linearized risk]\label{prop:conditionallinearized}
Under the static null, on the event
$\widetilde\lambda_r>\widetilde\lambda_{r+1}$, with
\[
H_t=\frac1W\bigl(r_t r_t^\top-r_{t-W}r_{t-W}^\top\bigr),
\]
one has
\[
\E\bigl[\norm{\mathcal L_{C_t}(H_t)}_F^2\bigm|\mathcal G_t\bigr]
=\frac4{W^2}
\sum_{i\le r<j}
\frac{a_i b_j+c_{ij}^2}{(\widetilde\lambda_i-\widetilde\lambda_j)^2}.
\]
\end{proposition}

\begin{proof}
For one boundary term $\Delta=xx^\top/W$ with $x\sim N(0,\Sigma)$, conditional
on $\mathcal G_t$,
\[
\norm{\mathcal L_{C_t}(\Delta)}_F^2
=\frac2{W^2}
\sum_{i\le r<j}
\frac{\abs{\widetilde u_j^\top x}^2\abs{\widetilde u_i^\top x}^2}
{(\widetilde\lambda_i-\widetilde\lambda_j)^2}.
\]
By Isserlis' formula \cite{Isserlis1918}, with $g_i=\widetilde u_i^\top x$ and
$g_j=\widetilde u_j^\top x$ jointly Gaussian,
\[
\E\bigl[g_i^2 g_j^2\bigm|\mathcal G_t\bigr]
=\E[g_i^2]\,\E[g_j^2]+2\bigl(\E[g_i g_j]\bigr)^2
=a_i b_j+2c_{ij}^2.
\]
Also, since $\E[\mathcal L_{C_t}(\Delta)\mid\mathcal G_t]=W^{-1}\mathcal L_{C_t}(\Sigma)$,
\[
\norm{\E[\mathcal L_{C_t}(\Delta)\mid\mathcal G_t]}_F^2
=\frac2{W^2}
\sum_{i\le r<j}
\frac{c_{ij}^2}{(\widetilde\lambda_i-\widetilde\lambda_j)^2}.
\]
Since the two boundary terms are conditionally i.i.d., the second moment of
their difference equals twice the conditional variance of one term, i.e.
\[
\E\bigl[\norm{\mathcal L_{C_t}(H_t)}_F^2\mid\mathcal G_t\bigr]
=2\Bigl(\E[\norm{\mathcal L_{C_t}(\Delta)}_F^2\mid\mathcal G_t]-\norm{\E[\mathcal L_{C_t}(\Delta)\mid\mathcal G_t]}_F^2\Bigr).
\]
Substituting the two displays gives
\[
\frac{2\cdot2}{W^2}\sum_{i\le r<j}\frac{(a_i b_j+2c_{ij}^2)-c_{ij}^2}{(\widetilde\lambda_i-\widetilde\lambda_j)^2}
=\frac4{W^2}\sum_{i\le r<j}\frac{a_i b_j+c_{ij}^2}{(\widetilde\lambda_i-\widetilde\lambda_j)^2}.
\]
\end{proof}

For later use, define the shared-core variance functional
\[
Q(C_t,\Sigma)
=
4\sum_{i\le r<j}
\frac{a_i b_j+c_{ij}^2}{(\widetilde\lambda_i-\widetilde\lambda_j)^2},
\]
with the measurable eigenbasis convention fixed above, so that
\[
\E\bigl[\norm{\mathcal L_{C_t}(H_t)}_F^2\bigm|\mathcal G_t\bigr]
=
\frac{Q(C_t,\Sigma)}{W^2}.
\]

\begin{lemma}[Shared-core linearized statistic is order $r$]\label{lem:sharedcorelinearized}
Assume the static null and work on $\mathcal E_t^{\mathrm{core}}$. If
$W\ge W_0$, so that $\alpha_W$ is bounded away from zero, then
\[
c r
\le
4\sum_{i\le r<j}
\frac{a_i b_j+c_{ij}^2}{(\widetilde\lambda_i-\widetilde\lambda_j)^2}
\le
C r.
\]
Consequently, for $W$ large enough,
\[
\E\bigl[
\norm{\mathcal L_{C_t}(H_t)}_F^2
\,\one_{\mathcal E_t^{\mathrm{core}}}
\bigr]
\asymp
\frac r{W^2}.
\]
\end{lemma}

\begin{proof}
On $\mathcal E_t^{\mathrm{core}}$,
\[
\widetilde\lambda_r-\widetilde\lambda_{r+1}
\ge
\alpha_W\delta_t-2\eta\,\alpha_W n
\ge
\frac{\alpha_W\delta_t}{2}
\asymp n.
\]
Hence for $i\le r<j$, $\abs{\widetilde\lambda_i-\widetilde\lambda_j}\gtrsim n$.
Also
\[
\widetilde\lambda_i
\le
\alpha_W\lambda_i(\Sigma)+\norm{C_t-\alpha_W\Sigma}_{\op}
\lesssim n
\qquad
(i\le r),
\]
and $C_t\succeq0$, so $0\le\widetilde\lambda_j\le\widetilde\lambda_i\lesssim n$
for $i\le r<j$. Thus
$(\widetilde\lambda_i-\widetilde\lambda_j)^2\asymp n^2$ for $i\le r<j$. By
Cauchy--Schwarz, $c_{ij}^2\le a_i b_j$, so
\[
\sum_{i\le r<j}(a_i b_j+c_{ij}^2)
\le
2\sum_{i\le r<j}a_i b_j
=2\Bigl(\sum_{i\le r}a_i\Bigr)\Bigl(\sum_{j>r}b_j\Bigr).
\]
Since $\Sigma$ has $r$ spikes of size $O(n)$ and bulk trace $O(n)$,
\[
\sum_{i\le r}a_i\le r\norm{\Sigma}_{\op}\lesssim rn,
\qquad
\sum_{j>r}b_j\le \tr\Sigma\lesssim n.
\]
Dividing by denominators of order $n^2$ gives the upper bound $Cr$.

For the lower bound, on $\mathcal E_t^{\mathrm{core}}$,
\[
\widetilde\lambda_i
\ge
\alpha_W\lambda_i(\Sigma)-\norm{C_t-\alpha_W\Sigma}_{\op}
\gtrsim n
\qquad
(i\le r).
\]
Moreover, since $\widetilde u_i$ is a unit eigenvector of $C_t$ with eigenvalue
$\widetilde\lambda_i$,
\[
\alpha_W a_i
=\widetilde u_i^\top(\alpha_W\Sigma)\widetilde u_i
\ge
\widetilde u_i^\top C_t\widetilde u_i-\norm{C_t-\alpha_W\Sigma}_{\op}
=\widetilde\lambda_i-\norm{C_t-\alpha_W\Sigma}_{\op}.
\]
Choosing $\eta$ sufficiently small gives $a_i\gtrsim n$ for $i\le r$. Also, by
the Ky Fan variational principle \cite[Ch.~III]{Bhatia1997},
\[
\sum_{j>r}b_j
=\tr\bigl((I-\widetilde P_t^C)\Sigma\bigr)
\ge
\sum_{k=r+1}^n\lambda_k(\Sigma)
\gtrsim n.
\]
Therefore, dropping $c_{ij}^2\ge0$ and bounding the denominators above by
$Cn^2$,
\[
4\sum_{i\le r<j}\frac{a_i b_j+c_{ij}^2}{(\widetilde\lambda_i-\widetilde\lambda_j)^2}
\ge
\frac{4}{Cn^2}\Bigl(\sum_{i\le r}a_i\Bigr)\Bigl(\sum_{j>r}b_j\Bigr)
\gtrsim
\frac{(rn)(n)}{n^2}
\asymp r.
\]
This is the lower bound $cr$. Proposition~\ref{prop:conditionallinearized} and
Lemma~\ref{lem:sharedcoreevent} then give the expectation statement.
\end{proof}

\begin{lemma}[Boundary tail at an arbitrary fixed threshold]\label{lem:tiptail}
Under the static null and Assumption~\ref{ass:bulkspike}, if
$r_\pm\sim N(0,\Sigma)$ independently, then for every fixed $\kappa>0$ there
exist constants $c_\kappa,C_\kappa,W_\kappa>0$, depending only on $\kappa$ and
on the spectral constants, such that
\[
\Prob\bigl\{\norm{r_\pm}_2^2>\kappa nW\bigr\}
\le
C_\kappa e^{-c_\kappa W}
\]
for all $W\ge W_\kappa$.
\end{lemma}

\begin{proof}
Write $\norm{r_\pm}_2^2=\sum_{k=1}^n\lambda_k Z_k^2$ with independent standard
normal $Z_k$. Assumption~\ref{ass:bulkspike} gives
\[
\tr\Sigma\le C_0 n,
\qquad
\tr(\Sigma^2)\le C_1 n^2,
\qquad
\norm{\Sigma}_{\op}\le C_2 n.
\]
For the fixed threshold $\kappa>0$, choose $W_\kappa\ge 2C_0/\kappa$, so that
for $W\ge W_\kappa$,
\[
\kappa nW-\tr\Sigma\ge \frac{\kappa}{2}nW.
\]
The Bernstein-type bound for weighted chi-square variables
\cite{LaurentMassart2000,HsuKakadeZhang2012} gives, for all $u>0$,
\[
\Prob\left\{\norm{r_\pm}_2^2-\tr\Sigma>u\right\}
\le
\exp\left[-c\min\left(\frac{u^2}{\tr(\Sigma^2)},\frac{u}{\norm{\Sigma}_{\op}}\right)\right].
\]
Substituting $u=(\kappa/2)nW$ yields, since $u^2/\tr(\Sigma^2)\gtrsim W^2$ and
$u/\norm{\Sigma}_{\op}\gtrsim W$,
\[
\Prob\bigl\{\norm{r_\pm}_2^2>\kappa nW\bigr\}
\le
\exp\left[-c\min\{c_\kappa W^2,c_\kappa W\}\right]
\le
C_\kappa e^{-c_\kappa W},
\]
after adjusting constants.
\end{proof}

\begin{lemma}[Fourth moment of the shared-core linearization]\label{lem:sharedcorefourth}
Assume the static null and work on $\mathcal E_t^{\mathrm{core}}$. Let
\[
L_t=\mathcal L_{C_t}\left(\frac1W(r_t r_t^\top-r_{t-W}r_{t-W}^\top)\right).
\]
Then
\[
\E\bigl[\norm{L_t}_F^4\bigm|\mathcal G_t\bigr]
\le
C\frac{r^2}{W^4}.
\]
\end{lemma}

\begin{proof}
Conditional on $\mathcal G_t$, $L_t$ is a Hilbert-valued Gaussian chaos of
order two, obtained as the difference of two independent copies passed through
the fixed linear map $\mathcal L_{C_t}$. Hypercontractivity for Gaussian
chaoses \cite[Ch.~5--6]{Janson1997} (see also \cite{Nelson1973}) gives
\[
\E\bigl[\norm{L_t}_F^4\bigm|\mathcal G_t\bigr]
\le
C\left(\E\bigl[\norm{L_t}_F^2\bigm|\mathcal G_t\bigr]\right)^2.
\]
By Proposition~\ref{prop:conditionallinearized} and the upper bound in
Lemma~\ref{lem:sharedcorelinearized}, the conditional second moment is at most
$Cr/W^2$ on $\mathcal E_t^{\mathrm{core}}$. Squaring gives the claim.
\end{proof}

\begin{lemma}[Shared-core derivative replacement on the good event]
\label{lem:corederivativeproved}
Assume the static null, Assumptions~\ref{ass:bulkspike}--\ref{ass:gap}, fixed
$r$, $W\to\infty$, and $W\ge C\log n$. Then
\[
\E\left[
\left|
Q(C_t,\Sigma)
-
\frac{2}{\alpha_W^2}\mathfrak d_{\mathrm{sub}}(\Sigma,r)
\right|
\one_{\mathcal E_t^{\mathrm{core}}}
\right]
\le
C\frac{\E\norm{C_t-\alpha_W\Sigma}_{\op}}{n}
=
O(W^{-1/2})
=o(1).
\]
\end{lemma}

\begin{proof}
By Lemma~\ref{lem:covopmoments} applied to
$S_t^{\mathrm{core}}=(W-1)^{-1}\sum_{s=t-W+1}^{t-1}r_s r_s^\top$,
\[
\E\norm{C_t-\alpha_W\Sigma}_{\op}
=\alpha_W\,\E\norm{S_t^{\mathrm{core}}-\Sigma}_{\op}
\le \alpha_W\bigl(\E\norm{S_t^{\mathrm{core}}-\Sigma}_{\op}^2\bigr)^{1/2}
\le
C\frac{n}{\sqrt W}.
\]

It remains to record the perturbative continuity of the variance functional.
Let $A_0=\alpha_W\Sigma$, and let $A$ be any symmetric matrix with
$\norm{A-A_0}_{\op}\le \eta\alpha_W n$ for $\eta$ sufficiently small. Then the
top-$r$ spectral cluster of $A$ is separated from its complement by a gap of
order $n$. Define
\[
Q(A,\Sigma)
=
\E\norm{\mathcal L_A(xx^\top-yy^\top)}_F^2,
\qquad
x,y\stackrel{\mathrm{iid}}{\sim}N(0,\Sigma),
\]
which for $A=C_t$ agrees with the coordinate formula above. By the
Riesz-projector representation of $\mathcal L_A$ \cite[Ch.~II]{Kato1995}, the
resolvent identity, and the fixed-rank bulk-plus-spikes moment bounds,
\[
\abs{Q(A,\Sigma)-Q(A_0,\Sigma)}
\le
C\frac{\norm{A-A_0}_{\op}}{n}
\]
uniformly on this neighborhood: the contour separating the top cluster has
length $O(n)$, all relevant resolvents are $O(n^{-1})$, the resolvent
difference contributes one factor $\norm{A-A_0}_{\op}$, and the Gaussian second
moments entering the off-block reduced-resolvent variance are $O(r)$.

At $A_0=\alpha_W\Sigma$ the eigenvectors agree with those of $\Sigma$ and the
eigenvalues are $\alpha_W\lambda_k$, so $a_i=\lambda_i$, $b_j=\lambda_j$,
$c_{ij}=0$, and every denominator is multiplied by $\alpha_W^2$. Hence
\[
Q(\alpha_W\Sigma,\Sigma)
=4\sum_{i\le r<j}\frac{\lambda_i\lambda_j}{\alpha_W^2(\lambda_i-\lambda_j)^2}
=\frac{4}{\alpha_W^2}\,\ssum(\Sigma,r)
=
\frac{2}{\alpha_W^2}\mathfrak d_{\mathrm{sub}}(\Sigma,r),
\]
using $\ssum=\tfrac12\mathfrak d_{\mathrm{sub}}$ from
Remark~\ref{rem:convention}. Combining this identity with the Lipschitz bound
and the operator-moment estimate proves the claim.
\end{proof}

\begin{theorem}[Adjacent-window drift benchmark under the static null]\label{thm:overlapbenchmark}
Assume the static null on $I_t\cup I_{t-1}$ and
Assumptions~\ref{ass:bulkspike}--\ref{ass:gap}, with $r$ fixed. There exist
constants $c,C>0$ and $W_0<\infty$, with $W_0$ depending on the fixed rank $r$
and the spectral constants, such that if
\[
W\ge C\log n
\qquad
\text{and}
\qquad
W\ge W_0,
\]
then
\[
c\,\frac r{W^2}
\le
\E\norm{\widehat P_t-\widehat P_{t-1}}_F^2
\le
C\,\frac r{W^2}+C r\, e^{-cW}.
\]
\end{theorem}

\begin{proof}
Let $\gamma_t^C=\widetilde\lambda_r-\widetilde\lambda_{r+1}$ be the top-$r$
shared-core eigengap when positive, and set
\[
\Delta_t^+=\frac1W r_t r_t^\top,
\qquad
\Delta_t^-=\frac1W r_{t-W} r_{t-W}^\top,
\qquad
H_t=\Delta_t^+-\Delta_t^-.
\]
Choose a small constant $\rho\in(0,c_0]$, with $c_0$ as in
Lemma~\ref{lem:projectorexpansion}, and define
\[
\mathcal H_t
=\mathcal E_t^{\mathrm{core}}
\cap
\left\{
\norm{\Delta_t^+}_{\op}\vee\norm{\Delta_t^-}_{\op}
\le
\rho\gamma_t^C
\right\}.
\]
On $\mathcal E_t^{\mathrm{core}}$ there is a deterministic constant $c_\gamma>0$
with $\gamma_t^C\ge c_\gamma n$. Choose the fixed threshold
$\kappa=\rho c_\gamma/2$. Then
\[
\frac{\norm{r_t}_2^2}{W}
\vee
\frac{\norm{r_{t-W}}_2^2}{W}
\le
\kappa n
\]
implies $\norm{\Delta_t^+}_{\op}\vee\norm{\Delta_t^-}_{\op}\le\rho\gamma_t^C$,
hence implies $\mathcal H_t$ on $\mathcal E_t^{\mathrm{core}}$.

The complement of $\mathcal H_t$ is contained in the union of
$(\mathcal E_t^{\mathrm{core}})^c$ and the event that at least one boundary
vector exceeds the fixed threshold $\kappa nW$.
Lemma~\ref{lem:sharedcoreevent} controls the first event and
Lemma~\ref{lem:tiptail}, applied with this $\kappa$, controls the second.
Therefore $\Prob(\mathcal H_t^c)\le Ce^{-cW}$.

On $\mathcal H_t$, Lemma~\ref{lem:projectorexpansion} applies at $C_t$ to both
$\Delta_t^+$ and $\Delta_t^-$, so
\[
\widehat P_t
=P(C_t)+\mathcal L_{C_t}(\Delta_t^+)+R_t^+,
\qquad
\widehat P_{t-1}
=P(C_t)+\mathcal L_{C_t}(\Delta_t^-)+R_t^-,
\]
with $\norm{R_t^\pm}_F\le C_r\norm{\Delta_t^\pm}_{\op}^2/(\gamma_t^C)^2$.
Subtracting,
\[
\widehat P_t-\widehat P_{t-1}
=\mathcal L_{C_t}(H_t)+R_t^\Delta,
\qquad
R_t^\Delta=R_t^+-R_t^-.
\]
On $\mathcal H_t$, since $\gamma_t^C\asymp n$,
\[
\E\bigl[\norm{R_t^\Delta}_F^2\,\one_{\mathcal H_t}\bigr]\le \frac{C_r}{W^4},
\]
using $\E\norm{r_\pm}_2^8\le Cn^4$ under the bulk-plus-spikes assumption and
fixed $r$.

Let $L_t=\mathcal L_{C_t}(H_t)$. By Lemma~\ref{lem:sharedcorelinearized},
$\E[\norm{L_t}_F^2\,\one_{\mathcal E_t^{\mathrm{core}}}]\asymp r/W^2$, and
Lemma~\ref{lem:sharedcorefourth} gives the matching fourth-moment bound.
Therefore, by Cauchy--Schwarz and $\Prob(\mathcal H_t^c)\le Ce^{-cW}$,
\[
\E\bigl[
\norm{L_t}_F^2
\,\one_{\mathcal E_t^{\mathrm{core}}\cap \mathcal H_t^c}
\bigr]
\le
\sqrt{\E[\norm{L_t}_F^4\,\one_{\mathcal E_t^{\mathrm{core}}}]}\,\sqrt{\Prob(\mathcal H_t^c)}
\le
C\frac r{W^2}e^{-cW/2},
\]
hence $\E[\norm{L_t}_F^2\,\one_{\mathcal H_t}]\asymp r/W^2$.

For the upper bound, use
$\norm{L_t+R_t^\Delta}_F^2\le 2\norm{L_t}_F^2+2\norm{R_t^\Delta}_F^2$ on
$\mathcal H_t$ and $\norm{\widehat P_t-\widehat P_{t-1}}_F^2\le 2r$ on
$\mathcal H_t^c$, giving
\[
\E\norm{\widehat P_t-\widehat P_{t-1}}_F^2
\le
C\frac r{W^2}+C r\, e^{-cW}.
\]
For the lower bound,
$\norm{L_t+R_t^\Delta}_F^2\ge\norm{L_t}_F^2-2\abs{\langle L_t,R_t^\Delta\rangle}$,
and by Cauchy--Schwarz
\[
2\,\E\bigl[\abs{\langle L_t,R_t^\Delta\rangle}\,\one_{\mathcal H_t}\bigr]
\le
2\sqrt{\E[\norm{L_t}_F^2\one_{\mathcal H_t}]}\,\sqrt{\E[\norm{R_t^\Delta}_F^2\one_{\mathcal H_t}]}
\le
C\frac{\sqrt r}{W^3}.
\]
Since $r$ is fixed, this is of lower order than $r/W^2$ for large $W$. Dropping
the nonnegative contribution off $\mathcal H_t$ and enlarging $W_0$ if necessary
gives $\E\norm{\widehat P_t-\widehat P_{t-1}}_F^2\ge c\,r/W^2$.
\end{proof}

\begin{theorem}[Constant-level adjacent-window drift expansion]
\label{thm:adjacentconstantproved}
Assume the static null on $I_t\cup I_{t-1}$,
Assumptions~\ref{ass:bulkspike}--\ref{ass:gap}, fixed $r$, $W\to\infty$, and
$W\ge C\log n$. Then
\[
\E\norm{\widehat P_t-\widehat P_{t-1}}_F^2
=
\frac{2}{\alpha_W^2}
\frac{\mathfrak d_{\mathrm{sub}}(\Sigma,r)}{W^2}
+
o(W^{-2}).
\]
Since $\alpha_W=(W-1)/W$ and
$\mathfrak d_{\mathrm{sub}}(\Sigma,r)=O(1)$ for fixed $r$, equivalently,
\[
\E\norm{\widehat P_t-\widehat P_{t-1}}_F^2
=
\frac{2\mathfrak d_{\mathrm{sub}}(\Sigma,r)}{W^2}
+
o(W^{-2}).
\]
\end{theorem}

\begin{proof}
Use the event $\mathcal H_t$ and decomposition from the proof of
Theorem~\ref{thm:overlapbenchmark}. On $\mathcal H_t$,
\[
\widehat P_t-\widehat P_{t-1}
=
L_t+R_t^\Delta,
\qquad
L_t=\mathcal L_{C_t}(H_t),
\]
with $\E[\norm{R_t^\Delta}_F^2\one_{\mathcal H_t}]\le C W^{-4}$ and
$\Prob(\mathcal H_t^c)\le Ce^{-cW}$; on the complement
$\norm{\widehat P_t-\widehat P_{t-1}}_F^2\le 2r$.

By Proposition~\ref{prop:conditionallinearized},
$\E[\norm{L_t}_F^2\mid\mathcal G_t]=Q(C_t,\Sigma)/W^2$ on the shared-core gap
event. Lemma~\ref{lem:corederivativeproved}, together with
Lemma~\ref{lem:sharedcoreevent}, gives
\[
\E\bigl[\norm{L_t}_F^2\one_{\mathcal E_t^{\mathrm{core}}}\bigr]
=
\frac{1}{W^2}
\left(
\frac{2}{\alpha_W^2}\mathfrak d_{\mathrm{sub}}(\Sigma,r)
+
o(1)
\right).
\]
The loss from further restricting to $\mathcal H_t$ is exponentially small: by
Lemma~\ref{lem:sharedcorefourth} and Cauchy--Schwarz,
\[
\E\bigl[\norm{L_t}_F^2\one_{\mathcal H_t^c\cap\mathcal E_t^{\mathrm{core}}}\bigr]
\le
C\frac r{W^2}e^{-cW/2}.
\]
Finally,
\[
\norm{L_t+R_t^\Delta}_F^2
=
\norm{L_t}_F^2
+
2\langle L_t,R_t^\Delta\rangle_F
+
\norm{R_t^\Delta}_F^2,
\]
where the remainder square is $O(W^{-4})$ and, by Cauchy--Schwarz,
\[
\abs{\E[\langle L_t,R_t^\Delta\rangle_F\one_{\mathcal H_t}]}
\le
\sqrt{\E[\norm{L_t}_F^2\one_{\mathcal H_t}]}
\sqrt{\E[\norm{R_t^\Delta}_F^2\one_{\mathcal H_t}]}
=
O(W^{-3}),
\]
since $r$ is fixed. Both terms, and the exponentially small off-event
contribution, are $o(W^{-2})$. This proves the stated expansion.
\end{proof}

\begin{remark}[Explicit leading constant and convention check]\label{rem:explicit}
In terms of the un-doubled cross-spectral sum
$\ssum(\Sigma,r)=\tfrac12\mathfrak d_{\mathrm{sub}}(\Sigma,r)$ of
Remark~\ref{rem:convention}, Theorem~\ref{thm:adjacentconstantproved} reads
\[
\E\norm{\widehat P_t-\widehat P_{t-1}}_F^2
=\frac{4}{W^2}\sum_{i\le r<j}\frac{\lambda_i\lambda_j}{(\lambda_i-\lambda_j)^2}
+o(W^{-2}),
\qquad
\E\,\tfrac12\norm{\widehat P_t-\widehat P_{t-1}}_F^2
=\frac{2\ssum(\Sigma,r)}{W^2}+o(W^{-2}).
\]
The leading constant coincides with the population boundary risk
$\E\norm{\mathcal L_\Sigma(H)}_F^2=2\mathfrak d_{\mathrm{sub}}/W^2$ of
Proposition~\ref{prop:dsub}, as it must: the random-core gap correction
$\alpha_W^{-2}\to1$ recovers the population linearization. Consistency with the
conditional functional of Proposition~\ref{prop:conditionallinearized} also
holds by the law of total expectation, since at the population spectrum
$Q(\alpha_W\Sigma,\Sigma)=2\mathfrak d_{\mathrm{sub}}/\alpha_W^2$, which divided
by $W^2$ matches the displayed constant.
\end{remark}

\section{Boundary-driven sharpness}

The adjacent-window lower bound requires a precise target. Under a known static
null the population instantaneous drift $P_t-P_{t-1}$ is zero and the estimator
$\widehat D\equiv0$ has zero risk. Under a fully nonstatic class, instantaneous
drift may be harder than $r/W^2$ because only one new observation is available
at each endpoint. We therefore state the sharp $r/W^2$ lower bound for the
overlap-driven rolling target
\[
D_t^{\mathrm{roll}}
=
P(\bar\Sigma_t)-P(\bar\Sigma_{t-1}),
\]
where $\bar\Sigma_t$ denotes the population window average, or equivalently for
local path models in which adjacent rolling drift is generated by the boundary
innovation.

\begin{proposition}[Boundary-driven rolling-drift lower bound]\label{prop:rollinglower}
Fix $r$ and suppose $n\ge2r+1$. There is a Gaussian adjacent-window subclass
satisfying the bulk-plus-spikes and eigengap conditions uniformly such that
\[
\inf_{\widehat D}\sup
\E\norm{\widehat D-D_t^{\mathrm{roll}}}_F^2
\ge
c\frac r{W^2}.
\]
\end{proposition}

\begin{proof}
Let $d=r(n-r)$ and parametrize local rotations by matrices
$\Theta\in\mathbb R^{(n-r)\times r}$ through
\[
U_\Theta=
\begin{pmatrix}
I_r\\
\Theta
\end{pmatrix}
(I_r+\Theta^\top\Theta)^{-1/2},
\qquad
P_\Theta=U_\Theta U_\Theta^\top,
\qquad
\Sigma_\Theta=I_n+nP_\Theta.
\]
For sufficiently small $\norm{\Theta}_{\op}$, this subclass has $r$ spike
eigenvalues equal to $n+1$, all remaining eigenvalues equal to $1$, and
eigengap $n$.

In the two adjacent windows, set the $W-1$ shared-core covariances equal to
$\Sigma_0=I_n+nP_0$, while the two boundary covariances are
$\Sigma_t=\Sigma_\Theta$ and $\Sigma_{t-W}=\Sigma_{-\Theta}$. Then
\[
\bar\Sigma_t=\Sigma_0+\frac1W(\Sigma_\Theta-\Sigma_0),
\qquad
\bar\Sigma_{t-1}=\Sigma_0+\frac1W(\Sigma_{-\Theta}-\Sigma_0).
\]
A first-order projector expansion at $\Sigma_0$ gives, uniformly for
sufficiently small $\norm{\Theta}_{\op}$,
\[
\norm{D_t^{\mathrm{roll}}(\Theta)-D_t^{\mathrm{roll}}(\Theta')}_F^2
\asymp
\frac1{W^2}\norm{\Theta-\Theta'}_F^2
\]
for points in the local chart (the spike size $n$ in
$\Sigma_{\pm\Theta}-\Sigma_0$ cancels against the order-$n$ eigengap in the
reduced resolvent, leaving an $O(1/W)$ map of $\Theta$). Only the two boundary
observations depend on $\Theta$, so the Kullback--Leibler divergence between
adjacent hypercube vertices is bounded by $C n h^2$, with no factor $W$.
Choosing $h^2=c_0/n$, Assouad's lemma \cite{Assouad1983,Yu1997,Tsybakov2009}
over the $d=r(n-r)$ tangent coordinates yields
\[
\inf_{\widehat D}\sup_\Theta
\E_\Theta\norm{\widehat D-D_t^{\mathrm{roll}}(\Theta)}_F^2
\ge
c\frac{d h^2}{W^2}
\asymp
\frac r{W^2}.
\]
\end{proof}

\begin{remark}[Instantaneous versus rolling drift]\label{rem:instantaneousvsrolling}
Proposition~\ref{prop:rollinglower} concerns $D_t^{\mathrm{roll}}$, not the
arbitrary instantaneous target $P_t-P_{t-1}$. A lower bound for instantaneous
drift requires an additional local path model linking $P_t-P_{t-1}$ to the
boundary difference seen by the two adjacent rolling windows. Without such a
link the problem can be either easier, as under a known static null, or harder,
as in unrestricted nonstationary models.
\end{remark}

\section{Discussion}

The central mechanism in this analysis is the shared-core decomposition of
adjacent rolling windows. It yields a null drift scale that is one order of
magnitude in $W$ smaller than the single-window projector error: the
single-window perturbation uses $W$ noisy observations, whereas the adjacent
difference uses only the entering and exiting observations after the common
core cancels. Under the bulk-plus-spikes scaling this produces the two scales
\[
\text{single-window projector scale } r/W,
\qquad
\text{adjacent overlap drift scale } r/W^2.
\]
The exact leading constant is not obtained by applying Davis--Kahan
\cite{DavisKahan1970} to two independent windows. It requires conditioning on
the shared core, expanding each endpoint around that same random core, and then
replacing the random-core variance functional by its population limit. Two
distinct calibrations follow from the same conditioning: the magnitude
benchmark $2\mathfrak d_{\mathrm{sub}}/W^2$ of
Theorem~\ref{thm:adjacentconstantproved}, which is asymptotic and carries the
factor-of-two convention of Remark~\ref{rem:convention}, and the exact
finite-sample sign symmetry of Theorem~\ref{thm:exactrandomization}, which
needs no constant at all and so remains calibrated for oriented contrasts even
when the magnitude scale is in doubt.

\begin{thebibliography}{99}

\bibitem{Assouad1983}
P.~Assouad.
\newblock Deux remarques sur l'estimation.
\newblock {\em C. R. Acad. Sci. Paris S\'er. I Math.}, 296(23):1021--1024, 1983.

\bibitem{AueEtAl2009}
A.~Aue, S.~H\"ormann, L.~Horv\'ath, and M.~Reimherr.
\newblock Break detection in the covariance structure of multivariate time
  series models.
\newblock {\em Ann. Statist.}, 37(6B):4046--4087, 2009.

\bibitem{BenaychNadakuditi2011}
F.~Benaych-Georges and R.~R. Nadakuditi.
\newblock The eigenvalues and eigenvectors of finite, low rank perturbations of
  large random matrices.
\newblock {\em Adv. Math.}, 227(1):494--521, 2011.

\bibitem{Bhatia1997}
R.~Bhatia.
\newblock {\em Matrix Analysis}, volume 169 of {\em Graduate Texts in
  Mathematics}.
\newblock Springer, New York, 1997.

\bibitem{CaiMaWu2013}
T.~T. Cai, Z.~Ma, and Y.~Wu.
\newblock Sparse {PCA}: optimal rates and adaptive estimation.
\newblock {\em Ann. Statist.}, 41(6):3074--3110, 2013.

\bibitem{CaiZhang2018}
T.~T. Cai and A.~Zhang.
\newblock Rate-optimal perturbation bounds for singular subspaces with
  applications to high-dimensional statistics.
\newblock {\em Ann. Statist.}, 46(1):60--89, 2018.

\bibitem{DavisKahan1970}
C.~Davis and W.~M. Kahan.
\newblock The rotation of eigenvectors by a perturbation. {III}.
\newblock {\em SIAM J. Numer. Anal.}, 7(1):1--46, 1970.

\bibitem{HsuKakadeZhang2012}
D.~Hsu, S.~M. Kakade, and T.~Zhang.
\newblock A tail inequality for quadratic forms of subgaussian random vectors.
\newblock {\em Electron. Commun. Probab.}, 17(52):1--6, 2012.

\bibitem{Isserlis1918}
L.~Isserlis.
\newblock On a formula for the product-moment coefficient of any order of a
  normal frequency distribution in any number of variables.
\newblock {\em Biometrika}, 12(1/2):134--139, 1918.

\bibitem{Janson1997}
S.~Janson.
\newblock {\em Gaussian Hilbert Spaces}, volume 129 of {\em Cambridge Tracts in
  Mathematics}.
\newblock Cambridge University Press, Cambridge, 1997.

\bibitem{Johnstone2001}
I.~M. Johnstone.
\newblock On the distribution of the largest eigenvalue in principal components
  analysis.
\newblock {\em Ann. Statist.}, 29(2):295--327, 2001.

\bibitem{Kato1995}
T.~Kato.
\newblock {\em Perturbation Theory for Linear Operators}.
\newblock Classics in Mathematics. Springer, Berlin, reprint of the 1980
  edition, 1995.

\bibitem{KoltchinskiiLounici2017}
V.~Koltchinskii and K.~Lounici.
\newblock Concentration inequalities and moment bounds for sample covariance
  operators.
\newblock {\em Bernoulli}, 23(1):110--133, 2017.

\bibitem{LaurentMassart2000}
B.~Laurent and P.~Massart.
\newblock Adaptive estimation of a quadratic functional by model selection.
\newblock {\em Ann. Statist.}, 28(5):1302--1338, 2000.

\bibitem{LehmannRomano2005}
E.~L. Lehmann and J.~P. Romano.
\newblock {\em Testing Statistical Hypotheses}.
\newblock Springer Texts in Statistics. Springer, New York, third edition,
  2005.

\bibitem{Nelson1973}
E.~Nelson.
\newblock The free {M}arkoff field.
\newblock {\em J. Funct. Anal.}, 12(2):211--227, 1973.

\bibitem{Paul2007}
D.~Paul.
\newblock Asymptotics of sample eigenstructure for a large dimensional spiked
  covariance model.
\newblock {\em Statist. Sinica}, 17(4):1617--1642, 2007.

\bibitem{StewartSun1990}
G.~W. Stewart and J.-G. Sun.
\newblock {\em Matrix Perturbation Theory}.
\newblock Academic Press, Boston, 1990.

\bibitem{Tsybakov2009}
A.~B. Tsybakov.
\newblock {\em Introduction to Nonparametric Estimation}.
\newblock Springer Series in Statistics. Springer, New York, 2009.

\bibitem{Vershynin2018}
R.~Vershynin.
\newblock {\em High-Dimensional Probability: An Introduction with Applications
  in Data Science}, volume 47 of {\em Cambridge Series in Statistical and
  Probabilistic Mathematics}.
\newblock Cambridge University Press, Cambridge, 2018.

\bibitem{VuLei2013}
V.~Q. Vu and J.~Lei.
\newblock Minimax sparse principal subspace estimation in high dimensions.
\newblock {\em Ann. Statist.}, 41(6):2905--2947, 2013.

\bibitem{WangYuRinaldo2021}
D.~Wang, Y.~Yu, and A.~Rinaldo.
\newblock Optimal covariance change point localization in high dimensions.
\newblock {\em Bernoulli}, 27(1):554--575, 2021.

\bibitem{Wedin1972}
P.-{\AA}. Wedin.
\newblock Perturbation bounds in connection with singular value decomposition.
\newblock {\em BIT}, 12(1):99--111, 1972.

\bibitem{Yu1997}
B.~Yu.
\newblock Assouad, {F}ano, and {L}e {C}am.
\newblock In {\em Festschrift for Lucien Le Cam}, pages 423--435. Springer, New
  York, 1997.

\end{thebibliography}

\end{document}